% ===========================================================================
\section{Reproducibility}
\label{app:repro}
% ===========================================================================

The inventory runs from a documented pipeline (\texttt{python3 run\_all.py}) that
executes the agricultural stages --- national inventory, provincial downscaling,
gridding, district downscaling, time series, Monte-Carlo, key-category, sensitivity,
and uncertainty-input analysis, completeness table, and a consolidated dashboard ---
and writes all tables and figures reported here. The exploratory land sector is an
\emph{optional} module (skipped under \texttt{--quick}; see Supplementary Note~S1)
and feeds none of the agricultural results. The agricultural core runs on the Python
standard library alone; the gridding and figures use an optional geospatial stack.
The shipped surfaces need nothing supplied: the FIRMS extraction and the MapSPAM
rice raster are both committed-or-held and hash-pinned, so ``one command'' means
one command (the default contract is set out below). Every emission factor, parameter, and
activity value is stored in a versioned data file, and a suite of more than 240 unit
tests checks the IPCC worked-example calculations, exact reconciliation at each scale
(national, provincial, gridded, district, and the district\,$\times$\,year panel),
emission-factor traceability, and the key-category, sensitivity, uncertainty, and
completeness outputs. The suite also covers the air-pollutant engine and its shared
nitrogen ledger --- ledger closure at every stage, the separation of the pre-field and
post-application volatilisation fractions, conservation between the \NHthree{}-N and
\NHthree{}+\NOx{} bases, the province and grid allocation of every pollutant, the
matched-year and spatial benchmarks, and the structural manure-management scenarios ---
together with a consistency check that every headline number and figure in this
manuscript still matches the output file it was generated from.

\paragraph{Code and data availability.} The full source code, data files and tests are
version-controlled; DOIs are minted at archival (none exists yet, and no identifier is
printed before it resolves), and the archived release corresponding to any submitted
version will let reviewers run the same code that produced these numbers. The archive contains the calculation engine for both
layers, the shared nitrogen ledger (\texttt{outputs/n\_ledger\_2024.csv}, every stage-flow
row as kilograms of nitrogen and as a share of what was excreted) and its diagnostic runs
(\texttt{outputs/n\_ledger\_stage\_comparison.csv}, the IPCC and EMEP ammonia-nitrogen
factors set against each other stage by stage), the versioned emission-factor
and activity files, the gridded (GeoTIFF, EPSG:4326) and district (CSV) products, the
1990--2024 time series, the uncertainty and benchmark outputs---including
\texttt{outputs/air\_\allowbreak uncertainty\_\allowbreak montecarlo.csv}, the
per-pollutant intervals with the Approach~1 comparison beside them---and the test suite; units,
coordinate reference system and time coverage are recorded with each product, and a run
manifest records the exact inputs behind the tables reported here. Reproduction: install the (optional) geospatial stack and
run \texttt{python3 run\_all.py}. The shipped surfaces reproduce as-is with no
user-supplied input; the contract is stated in full in \ref{app:data}. Expected outputs are the CSV tables and PNG figures
reproduced here (written under \texttt{outputs/}). The deposited capsule carries its own
single command, \texttt{sh capsule/run\_tests.sh}, which reruns the agriculture suite in a
fresh interpreter with nothing installed; in the full repository
\texttt{python3 -m pytest tests} runs the
test suite, as does \texttt{python3 -m unittest discover tests}, which is what continuous
integration uses. On the released commit both pass in full; on any other commit the
release-integrity tests fail by design: they compare the
committed release artefacts---manifest, cover, checksums, PDFs---against the working tree,
and they can only pass on the exact commit the artefacts were cut from. They are not a failed
reproduction, and a reader who sees them should not read them as one. The code is released under the MIT licence and the derived data under
CC-BY-4.0; third-party inputs retain their original licences.

\section{Data acquisition}
\label{app:data}

THE REPRODUCTION CONTRACT, stated once. Livestock density (\GLW{}, 2020) is retrieved
directly from the FAO open cloud store. The FIRMS VIIRS extraction is \emph{committed}
and hash-pinned in \texttt{data/gis/manual/}, so no access key is needed to reproduce
(a key is needed only to RE-FETCH, which would create a new extraction). MapSPAM~2020
is \emph{held} (V2r2, hash-pinned in \texttt{RICE\_SOURCES.json}) and its former
cropland-fallback class is deleted from the allow-list, so a checkout missing the
raster refuses the build rather than silently substituting. The DEFAULT reproducible
surface is therefore: MapSPAM for rice, frozen FIRMS for burning, and nothing to
supply. Re-fetching FIRMS or swapping the rice surface (the own-year GCD-Rice 2023
classification ships as the priced alternative) changes the surface and is an
\emph{upgrade}: it creates a new configuration and a new manifest, and its outputs
are not the ones reported here. The
pipeline auto-detects supplied files and logs which proxy each category used.

\section{IPCC common reporting table}
\label{app:crt}

Table~\ref{tab:crt} gives a Common-Reporting-Table view of the agriculture sector
for the base year (1990) and the latest year (2024): the underlying activity data,
the implied emission factor, and the resulting \COe{}. The principal emission and
carbon factors themselves are listed in the body (Table~\ref{tab:ef-full}).

\begin{table}[!htbp]
\centering
\footnotesize
\setlength{\tabcolsep}{3pt}
\begin{tabular}{@{}p{0.22\linewidth}lp{0.30\linewidth}rr@{}}
\toprule
\textbf{Code / category} & \textbf{Gas} & \textbf{Activity (2024)}
& \textbf{1990} & \textbf{2024} \\
\midrule
3.A Enteric ferment. & \CHfour{} & 245.52 M ruminant head & 54.5 & 151.7 \\
3.B Manure mgmt.\ & \CHfour{} & 245.52 M ruminant head$^{a}$ & 4.3 & 13.2 \\
3.B Manure mgmt.\ & \NtwoO{} & 7.78 M t N excreted$^{a}$ & 2.4 & 8.0 \\
3.C Rice cultivation & \CHfour{} & 3.64 M ha paddy & 9.6 & 16.5 \\
3.D Managed soils & \NtwoO{} & F$_{SN}$ 3.78, F$_{ON}$ 2.34, F$_{PRP}$ 3.18, F$_{CR}$ 0.21 M t N$^{d}$ & 10.3 & 28.5 \\
3.F Field burning & \CHfour{},\NtwoO{} & 7.11 M t d.m.\ combusted$^{c}$ & 0.3 & 0.7 \\
3.H Urea application & \COtwo{} & 6.58 M t urea$^{b}$ & 1.9 & 4.8 \\
\midrule
\textbf{Agriculture total} & & & \textbf{83.2} & \textbf{223.5} \\
\bottomrule
\end{tabular}
\caption{IPCC CRT-style summary by category and gas, 1990 and 2024 (\COe{} in Mt,
AR5). $^{a}$Manure CH\textsubscript{4} is driven by the ruminant herd and its
N\textsubscript{2}O by the nitrogen excreted by that herd; both are independent of
the synthetic-fertiliser series. Poultry adds to the head-count but little emission. $^{b}$Urea mass is NFDC's printed national
offtake for the fiscal year; earlier years, which have no published product split, keep the
derivation from the national synthetic-N series, and each year states which it used.
$^{c}$The residue activity is dry matter actually
\emph{combusted}: production $\times$ residue ratio $\times$ burned fraction gives
\SI{8.4}{}~Mt of dry matter reaching the fire, of which the combustion factor
$C_f=0.85$ leaves \SI{7.1}{}~Mt burnt. Emission factors are applied to the latter. The full
destination accounting is published in
\texttt{outputs/agriculture\_crop\_residue\_ledger\_2024.csv}.
Liming (3.G): not estimated, assumed negligible on
Pakistan's alkaline, calcareous soils pending national lime-use data.
$^{d}$The four nitrogen inputs of Eq.~11.1, from the shared ledger: synthetic
fertiliser, applied organic manure, pasture deposition and crop residue. They are not
summed here because the pasture term carries its own factor (EF$_{3\mathrm{PRP}}$).}
\label{tab:crt}
\end{table}

\section{Completeness and QA/QC}
\label{app:completeness}

Table~\ref{tab:completeness} gives the notation-key coverage of the agriculture
categories. All major agriculture source categories are represented, with the
not-estimated (NE) gaps shown explicitly: soil-organic nitrogen mineralisation
($F_{SOM}$), which requires field data; and liming, set to NE (assumed negligible on
Pakistan's alkaline, calcareous soils pending national lime-use data) rather than
assumed not to occur. Dung-cake combustion is the one cross-sector transfer to
Energy. Land-sector completeness (forest, cropland/grassland SOC, wetlands,
harvested wood products) is documented in Supplementary Note~S1. Air-pollutant
co-emissions (\NHthree{}, \NOx{}, PM$_{2.5}$, CO, NMVOC, BC) are estimated for the
relevant categories in Section~\ref{sec:airpollutants} and reported in native mass,
outside this \COe{} completeness table; agricultural SO$_2$ is estimated but small
($\sim$2.6~kt, from residue burning only). Quality control is automated (source
traceability, unit checks, exact reconciliation, benchmark bands); quality
assurance is the independent FAOSTAT and EDGAR comparison; and any method or factor
change --- such as the harmonisation of the volatilisation factor $EF_4$ to the
dry-climate value across the manure and soils terms --- triggers a full
1990--2024 recalculation to preserve time-series consistency.

\begin{table}[!htbp]
\centering
\footnotesize
\begin{tabular}{@{}llc@{}}
\toprule
\textbf{Category} & \textbf{Gas} & \textbf{Key} \\
\midrule
3.A Enteric fermentation        & \CHfour{}          & E \\
3.B Manure management           & \CHfour{}, \NtwoO{} & E \\
3.C Rice cultivation            & \CHfour{}          & E \\
3.D Agricultural soils          & \NtwoO{}           & E \\
\quad SOC mineralisation ($F_{SOM}$) & \NtwoO{}      & NE \\
3.E Savanna burning             & ---                & NO \\
3.F Field burning of residues   & \CHfour{}, \NtwoO{} & E \\
3.G Liming                      & \COtwo{}           & NE \\
3.H Urea application            & \COtwo{}           & E \\
Dung-cake combustion            & multi              & IE \\
\bottomrule
\end{tabular}
\caption{Completeness and notation keys for the agriculture sector (IPCC Vol.~1
Ch.~8). E~estimated, NE~not estimated, NO~not occurring, IE~included elsewhere
(Energy). Land-sector categories are in Supplementary Note~S1.}
\label{tab:completeness}
\end{table}

\section{Uncertainty inputs}
\label{app:unc}

The uncertainty analysis (Section~\ref{sec:methods}) is reproducible from the inputs
below. Each category's combined uncertainty is a two-stage quadrature: every emission
component combines an activity-data and an emission-factor uncertainty (Eq.~3.1), and
the components are aggregated to the category (Eq.~3.2). The full underlying-input
table --- per component, the activity-data uncertainty, the factor uncertainty, their
sources, and which inputs are shared across categories --- is given in Supplementary
Note~S3; the factor uncertainties also appear as the $\pm$\% column of
Table~\ref{tab:ef-full}. Table~\ref{tab:unc} gives the resulting per-category
distribution used by the ANALYTIC comparison alone---the Monte-Carlo samples the engine's
components, not these fitted category distributions: the point estimate $E_c$, the combined
95\% half-width, and the mean-preserving lognormal parameters
$\sigma_{\ln}=(\text{half-width}/100)/1.96$ and
$\mu_{\ln}=\ln E_c-\sigma_{\ln}^2/2$. Correlation is modelled in one specific place and
not elsewhere, and the distinction matters: components that are separate rows only
because they follow different spatial proxies --- the managed-soils indirect
volatilisation and leaching streams --- share a single emission factor ($EF_4$ and
$EF_5$ respectively), so they are treated as perfectly correlated, summed linearly
before the quadrature step and given one shared deviate per group in the
Monte-Carlo. Treating them as independent would shrink the reported band purely
because the rows were cut more finely. All other within- and cross-category
correlations (shared activity data, common factor families across categories) are
\emph{not} modelled as shared deviates \emph{in this sector-scoped stream}, which samples
component totals; the national run's parameter graph does carry two of them, and
Supplementary Note~S3 reports their effect. Approach~1 combines these analytically; the
Monte-Carlo samples the \emph{components} themselves---each row drawn
mean-preserving lognormal at its declared width, with the correlated soils groups
receiving one shared deviate per draw---and the category and sector totals are the
sums of those sampled rows in every realisation.

\begin{table}[!htbp]
\centering
\footnotesize
\begin{tabular}{@{}lcccc@{}}
\toprule
\textbf{Category} & \textbf{$E_c$ (Mt)} & \textbf{$\pm$\% (95\%)} & \textbf{$\mu_{\ln}$} & \textbf{$\sigma_{\ln}$} \\
\midrule
Enteric fermentation & 151.7 & 25.4 &  5.013 & 0.129 \\
Manure management    &  21.2 & 32.9 &  3.042 & 0.167 \\
Managed soils        &  28.5 & 47.1 &  3.323 & 0.237 \\
Rice cultivation     &  16.6 & 60.9 &  2.760 & 0.304 \\
Field burning        &   0.7 & 44.9 & $-0.426$ & 0.226 \\
Urea application     &   4.8 & 11.2 &  1.571 & 0.057 \\
\midrule
\textbf{Total}       & \textbf{223.4} & \textbf{19.1} & --- & --- \\
\bottomrule
\end{tabular}
\caption{Derived per-category uncertainty parameters. $\mu_{\ln},\sigma_{\ln}$ are
the mean-preserving lognormal parameters implied by each category's 95\% half-width, and
they describe an ANALYTIC comparison, not the simulation: the Monte-Carlo samples the
engine's \emph{components} (with shared deviates inside declared correlation groups) and
obtains each category as the sum of its sampled rows. The national total is likewise a
sum, never a sampled input distribution, so no $\mu_{\ln}/
\sigma_{\ln}$ is given for it. Underlying activity/factor inputs and sources are in
Supplementary Note~S3. Generated by \texttt{src/pkinv/eval/uncertainty\_inputs.py}.}
\label{tab:unc}
\end{table}
